\documentclass[10pt,a4paper]{amsart}
\usepackage[margin=19mm]{geometry}
\usepackage{amsmath,amssymb,mathtools}
\usepackage[scr=rsfs,cal=euler]{mathalfa}
\usepackage{xfrac}
\usepackage[unicode,hidelinks,bookmarks=false]{hyperref}
\newcommand{\ie}{i.e.\ }
\newcommand{\cf}{cf.\ }
\newcommand{\resp}{respectively\ }
\newcommand{\Subsection}{\subsection}
\providecommand{\nobreakdash}{\nobreak\hbox{-}}
\setlength{\emergencystretch}{2em}
\multlinegap0pt
\usepackage{enumitem}
\usepackage{tikz}
\usetikzlibrary{decorations.pathreplacing}
\usepackage{subcaption}
\setlist[enumerate]{leftmargin=*, label=(\alph*)}
\setlist[itemize]{leftmargin=*}
\theoremstyle{plain}
\newtheorem{theorem}{Theorem}
\newtheorem{corollary}[theorem]{Corollary}
\newtheorem{lemma}[theorem]{Lemma}
\newcommand{\ceiling}[1]{\lceil #1 \rceil}
\newcommand{\floor}[1]{\lfloor #1 \rfloor}
\setcounter{theorem}{5}
\setcounter{equation}{5}
\title{The error term in the truncated Perron formula for the logarithm of an $L$-function: the sharp eight-prime sieve bound (Lemma 6)}
\author{Stephan Ramon Garcia, Jeffrey Lagarias, Ethan Simpson Lee}
\date{Source: arXiv:2206.01391v3; \url{https://arxiv.org/abs/2206.01391v3}}
\begin{document}
\maketitle
\noindent\textit{Editorial scope.} Counted unit: the article's sharp eight-prime sieve bound, printed as Lemma 6 (source0.tex lines 295--314, source label \texttt{Lemma:Sieve30}), delivered together with the article's own proof of it (source0.tex lines 336--369, the proof block that follows the lemma and whose parts (a) and (b) prove exactly the two displayed bounds of the lemma). The proof is the article's own text line for line, with no line added, dropped, reordered or re-flowed: part (a) reduces the first bound to $x \leq 9.5$ and then to the extremal pattern of eight primes in $(\frac{x}{2},x)$ modulo $30$, and part (b) does the same for the eight smallest primes in $(x,\frac{3x}{2})$; both retain the authors' stated finite computations exactly as printed, together with the commented-out prime-octuplet line inside part (a). No mathematics is written, repaired or re-derived, and the printed slips are kept literally.
Required context retained uncounted: the article's sentence introducing the lemma (line 272); and the article's figure carrying the labels \texttt{Figure:C1} and \texttt{Figure:C2} that the two parts of the statement cross-refer to (lines 316--333). That figure is retained with its subcaptions and labels; its two plot-inclusion commands are removed, and that removal is the only difference between the retained context text and the source lines, declared in the fidelity receipt of that body (exactly two removed strings, and nothing else). The two plot files of the e-print are not redistributed by this package.
Editorial formatting only, in addition: an AMS \texttt{amsart} wrapper that restates the article's own declarations that the retained lines use (the \texttt{theorem}, \texttt{corollary} and \texttt{lemma} environments sharing one counter exactly as the article declares them, the \texttt{enumerate} list style with (a), (b) labels, the macros \texttt{\textbackslash ceiling} and \texttt{\textbackslash floor} with the article's own definitions, and the packages \texttt{enumitem}, \texttt{tikz} with \texttt{decorations.pathreplacing}, and \texttt{subcaption} that the retained figure needs); three counter settings so that the retained lemma prints as Lemma 6, its two numbered displays print as (6) and (7) and the retained figure prints as Figure 2, their numbers in the article. The article is typeset in the publisher's own class with its own fonts, so the visual presentation differs although the mathematical text does not.
Bibliography: no retained line of this item cites anything, so the item carries no bibliography and no entry of the article's own bibliography is reproduced.
The next lemma is needed later to handle a few exceptional primes.

\begin{lemma}\label{Lemma:Sieve30}
    Let $x>1$ be a half integer and let $C = \frac{1284699552}{444215525}= 2.89206\ldots$.  
    \begin{enumerate}
        \item Let $p_{-8} < p_{-7} < \cdots < p_{-1} < x$ denote the largest eight primes (if they exist) in the interval $(\frac{x}{2},x)$.  
            We have the sharp bound\label{p:C1}
            \begin{equation}\label{eq:Mod30a}
                F_1(x) = \sum_{1 \leq n \leq 8}  \frac{1}{  x-p_{-n} } \leq  C;
            \end{equation}
            see Figure \ref{Figure:C1}.
            The corresponding summand in \eqref{eq:Mod30a} is zero if $p_{-n}$ does not exist.  
        
        \item Let $x<p_1 <p_2<\cdots < p_8$ denote the smallest eight primes (if they exist)
            in the interval $(x,\frac{3x}{2})$.  We have the sharp bound\label{p:C2}
            \begin{equation}\label{eq:Mod30b}
                F_2(x) = \sum_{1 \leq n \leq 8} \frac{1}{p_n -x} \leq  C;
            \end{equation}
		see Figure \ref{Figure:C2}.
            The corresponding summand in \eqref{eq:Mod30b} is zero if $p_{n}$ does not exist.  
    \end{enumerate}
\end{lemma}

\setcounter{figure}{1}
\begin{figure}
    \centering
    \begin{subfigure}[t]{0.47\textwidth}
    \centering

    \caption{$F_1(x)$ is bounded above by $C$ (red).}
    \label{Figure:C1}
    \end{subfigure}
    \quad
    \begin{subfigure}[t]{0.47\textwidth}
    \centering

    \caption{$F_2(x)$ is bounded above by $C$ (red).}
    \label{Figure:C2}
    \end{subfigure}
    \caption{The functions $F_1(x)$ and $F_2(x)$ behave erratically.}
    \label{Figure:C1C2}
\end{figure}

\begin{proof}
    (a) If $x \geq 10.5$, then $2,3,5 \notin (\frac{x}{2},x)$.  Computation confirms that
    $F_1(x) \leq F_1(3.5) = \frac{8}{3} =2.66\ldots$ for $x \leq 9.5$.  Let $x \geq 10.5$.
    Then any prime in $(\frac{x}{2},x)$ is 
    congruent to one of $1, 7, 11, 13, 17, 19, 23, 29 \pmod{30}$.  
    There are finitely many patterns modulo $30$ that the $p_{-8},p_{-7},\ldots,p_{-1}$ may assume.
    Among these, computation confirms that $F_1(x)$ is maximized if 
    \begin{align*}
        p_{-1} &= \floor{x} \equiv 19\pmod{30},  & p_{-5} &= \floor{x}-12 \equiv 7\pmod{30},\\
        p_{-2} &= \floor{x}-2 \equiv 17\pmod{30},  & p_{-6} &= \floor{x}-18 \equiv 1\pmod{30},\\
        p_{-3} &= \floor{x}-6 \equiv 13\pmod{30},  & p_{-7} &= \floor{x}-20 \equiv 29\pmod{30},\\
        p_{-4} &= \floor{x}-8 \equiv 11\pmod{30},  & p_{-8} &= \floor{x}-26 \equiv 23\pmod{30},
    \end{align*}
    which yields the desired upper bound $C$.
    This prime pattern first occurs for $x = 88819.5$;
%     which yields the prime octuplet
%    $88793$, $88799$, $88801$, $88807$, $88811$, $88813$, $88817$, $88819$;
    see \url{https://oeis.org/A022013}.
        
    \medskip\noindent(b) If $x \geq 5.5$, then $2,3,5 \notin (x, \frac{3x}{2})$.  Observe that
    $F_2(x) \leq 2$ for $x \leq 4.5$ (attained at $x=1.5, 2.5, 4.5$).  Let $x \geq 5.5$.
    As in (a), any prime in $(x,\frac{3x}{2})$ is 
    congruent to one of $1, 7, 11, 13, 17, 19, 23, 29 \pmod{30}$.  
    Then $F_2(x)$ is maximized if 
    \begin{align*}
        p_{1} &= \ceiling{x} \equiv 11\pmod{30},  & p_{5} &= \ceiling{x}+12 \equiv 23\pmod{30},\\
        p_{2} &= \ceiling{x}+2 \equiv 13\pmod{30},  & p_{6} &= \ceiling{x}+18 \equiv 29\pmod{30},\\
        p_{3} &= \ceiling{x}+6 \equiv 17\pmod{30},  & p_{7} &= \ceiling{x}+20 \equiv 1\pmod{30}, \\
        p_{4} &= \ceiling{x}+8 \equiv 19\pmod{30}, & p_8 &= \ceiling{x} + 26 \equiv 7 \pmod{30}.
    \end{align*}
    which yields the desired upper bound $C$.
    Although this prime pattern occurs for $x=10.5$, not all eight primes lie in $(x,\frac{3}{2}x)$.    
    The first admissible value is $x = 15760090.5$; see \url{https://oeis.org/A022011}.
\end{proof}

\end{document}
